\section{Two Lines, Swapped: A Case Study in Not Trusting a Pin Map}
\label{sec:i2c-swapped-lines}

A six-axis IMU was supposed to be on this board at address \codeinline{0x6A}. The vendor's
documentation said so, the vendor's example code said so, and the board's own pin map said
which two pins carried the bus. Every scan of those pins returned an empty bus: not a
mismatch, not garbage — nothing at all, as though no device existed.

It existed. The two lines were the other way round.

\subsection{What the Documentation Said}

Three sources, all consistent and all wrong in the same direction:

\begin{itemize}
    \item The vendor's product page: the IMU answers at \codeinline{0x6A}.
    \item The vendor's example: \codeinline{LSM6DS3 myIMU(I2C\_MODE, 0x6A)}, and a library
          header that selects \codeinline{Wire1} rather than the default bus.
    \item The board support package's pin map: \codeinline{SDA1 = PB3},
          \codeinline{SCL1 = PB2}.
\end{itemize}

The third line is the one that matters, and it is the one that is wrong. On this board the
data line is \textbf{PB2} and the clock is \textbf{PB3}.

\begin{warningbox}
    \textbf{Why a swapped pair produces silence rather than an error.} Both lines idle high
    through their pull-ups. If the master drives what it believes to be the clock while the
    device is watching it as data, the START condition — SDA falling while SCL is high —
    never occurs. The device never wakes up, never acknowledges, and never reports anything,
    because no transaction ever began. Every failure mode of a wrong pin assignment looks
    like an empty bus, not like a bad byte.
\end{warningbox}

\subsection{How It Was Found}

The useful part of this story is not the answer; it is the method, which is reusable on any
board with an unknown bus.

\paragraph{Step one: prove the pins are a bus at all.}
Before talking to anything, configure every GPIO as an input with no internal pull, and read
the port. Two pins that read high with no internal help are held up by resistances on the
board — which is what an I2C bus looks like, and nothing else on a general GPIO header does.

\begin{minted}{text}
PORTB_DIN = 0x0C    ->  bit 2 (PB2) and bit 3 (PB3) are both high, no internal pull
\end{minted}

\begin{notebox}
    The first version of this probe was useless: it configured the pins as \emph{outputs}
    (\codeinline{mode 0}), which disables the input buffer, so the result carried no
    information at all. \codeinline{mode 1} — input, no pull — is the measurement.
\end{notebox}

\paragraph{Step two: stop assuming which pin is which.}
Instead of testing one assignment, test the matrix. Every plausible pin pair, in both
orders, against both addresses the device could answer on:

\begin{minted}{text}
8 pin pairs x 2 orders x 2 addresses = 32 probes
result: sda=p1.2 scl=p1.3 addr=0x6a id=0x6a
\end{minted}

That single line is the whole fix. The device had been there the entire time, on the pair
the documentation listed — with the two wires the other way round.

\paragraph{Step three: make the driver's configuration a value, not a hardcoded write.}
The first attempt after the fix still failed, for a second and much more embarrassing reason:
the pin-abstraction's write function always addressed the same port, so changing ``which
bus'' changed nothing about the electrical output while appearing to work. The lesson is
small and general:

\begin{minted}{rust}
// Wrong: the selector looked configurable but every branch wrote PORTB.
fn set(&mut self, pin: u8, high: bool) {
    self.portb_dout().write(|w| unsafe { w.dout(if high { 1 << pin } else { 0 }) });
}

// Right: the port is part of the state, so choosing a bus chooses the registers.
fn write(&mut self, high: bool) {
    match self.pins {
        Pins::Base => self.portc_dout().write(|w| unsafe { w.dout(high as u32) }),
        Pins::Imu  => self.portb_dout().write(|w| unsafe { w.dout(high as u32) }),
    }
}
\end{minted}

\begin{remark}
    \textbf{The pattern to notice.} A parameter that does not reach the hardware is worse than
    a missing parameter: it produces a plausible-looking experiment whose result is meaningless,
    and you will believe the result.
\end{remark}

\subsection{Verification}

The fix was accepted only when two independent implementations agreed:

\begin{table}[H]
\centering
\begin{tabular}{lll}
\hline
\textbf{Implementation} & \textbf{Reading} & \textbf{Expected} \\
\hline
PAC example   & \codeinline{id=6a a=985,81,139 mg} & WHO\_AM\_I = 0x6A, $|a| \approx 1$ g \\
HAL example   & \codeinline{id=6a a=970,79,142 mg}  & same device, same registers \\
\hline
\end{tabular}
\caption{Two layers, one device: the numbers differ only by sensor noise and the physical
orientation at the moment of the read.}
\label{tab:i2c-imu-parity}
\end{table}

A magnitude of about 1000\,mg is not a detail to skip: a three-axis accelerometer at rest
measures gravity, so the vector length is a free self-check on both the bus and the
scale-factor arithmetic. A dead sensor reports zeros; a mis-scaled one reports something
that is not 1\,g.

\subsection{What This Teaches}

\begin{keytakeaways}
    \item \textbf{A pin map is a claim, not a measurement.} Check the electrical behaviour of
          the pins before believing a table in a repository.
    \item \textbf{Test the matrix, not the hypothesis.} Eight pairs in two orders took less
          time than a day of reasoning about which document was right.
    \item \textbf{Silence is informative.} An empty bus narrows the problem to ``no
          transaction started'', which points at wiring, mode or timing — never at the
          protocol layer.
    \item \textbf{Verify with a physical invariant.} Gravity, a chip identifier, a checksum:
          something that would be wrong if the bus were returning plausible garbage.
\end{keytakeaways}
